\documentclass[10pt,letterpaper]{article}

\usepackage[margin=0.52in]{geometry}
\usepackage[T1]{fontenc}
\usepackage[utf8]{inputenc}
\usepackage[default,type1]{sourcesanspro}
\usepackage[final]{microtype}
\usepackage[explicit]{titlesec}
\usepackage{enumitem}
\usepackage{tabularx}
\usepackage{array}
\usepackage{ragged2e}
\usepackage{iftex}
\usepackage[hidelinks]{hyperref}

\ifPDFTeX
  \input{glyphtounicode}
  \pdfgentounicode=1
\fi

\hypersetup{
  pdftitle={Connor Savugot Resume},
  pdfauthor={Connor Savugot},
  pdfsubject={Low-Level Embedded Firmware and Power/Control Systems Engineering},
  pdfkeywords={low-level embedded firmware, embedded C, C++, bare metal, FreeRTOS, peripheral bring-up, firmware update, dsPIC33CK, PIC32, TMS320 DSP, TI HAL, power conversion, CAN 2.0, CAN FD, J1939, DroneCAN, SPI, I2C, UART, VPX, IPMI, Embedded Linux, Yocto, Arago}
}

\pagestyle{empty}
\setcounter{secnumdepth}{0}
\setlength{\parindent}{0pt}
\setlength{\parskip}{0pt}
\hfuzz=0.4pt

\titleformat{\section}
  {\fontsize{11.3}{11.9}\selectfont\bfseries\uppercase}
  {}{0pt}{#1\hspace{0.12cm}\titlerule[0.65pt]}
\titlespacing{\section}{0pt}{0.15cm}{0.055cm}

\newlist{resumebullets}{itemize}{1}
\setlist[resumebullets]{
  label=\textbullet,
  leftmargin=0.42cm,
  topsep=0.005cm,
  itemsep=0.005cm,
  parsep=0pt,
  partopsep=0pt
}

\newcommand{\companyheader}[1]{%
  {\fontsize{10.4}{10.8}\selectfont\bfseries #1}\par\vspace{-0.03cm}
}

\newcommand{\resumeentry}[2]{%
  \begin{tabularx}{\textwidth}{@{}X>{\raggedleft\arraybackslash}p{2.40in}@{}}
    \textbf{#1} & #2
  \end{tabularx}\vspace{-0.075cm}
}

\newcommand{\roleentry}[2]{%
  \begin{tabularx}{\textwidth}{@{}X>{\raggedleft\arraybackslash}p{2.40in}@{}}
    \textbf{#1} & #2
  \end{tabularx}\vspace{-0.075cm}
}

\newcommand{\descriptorentry}[1]{%
  \begin{tabularx}{\textwidth}{@{}X>{\raggedleft\arraybackslash}p{2.40in}@{}}
    \textit{#1} &
  \end{tabularx}\vspace{-0.075cm}
}

\begin{document}
\fontsize{9.8}{10.72}\selectfont
\setlength{\RaggedRightRightskip}{0pt plus 5em}
\RaggedRight
\emergencystretch=1em
\hyphenpenalty=10000
\exhyphenpenalty=10000

\begin{center}
  {\fontsize{21.5}{22.2}\selectfont\bfseries Connor Savugot}\\[1.5pt]
  {\bfseries B.S. in Computer Engineering | North Carolina State University | Graduated May 2026}\\[1.5pt]
  {\fontsize{9.6}{10.2}\selectfont
    \href{mailto:consav04@gmail.com}{consav04@gmail.com}
    \enspace|\enspace \href{tel:+18285419487}{+1-828-541-9487}
    \enspace|\enspace Murphy, NC
    \enspace|\enspace \href{https://linkedin.com/in/connorsavugot}{linkedin.com/in/connorsavugot}
    \enspace|\enspace \href{https://github.com/Clem085}{github.com/Clem085}}
\end{center}
\vspace{-0.18cm}

\section{Technical Profile}
Low-level embedded firmware engineer specializing in real-time power, control, and communications software for microcontrollers and DSPs. Develops C/C++ firmware from peripheral bring-up and hardware debugging through firmware updates, system integration, and validation, with Embedded Linux as a supporting integration platform.

\section{Technical Skills}
\textbf{Low-Level Firmware:} Embedded C/C++, bare-metal and FreeRTOS development; peripheral bring-up, boot and update control\\
\textbf{MCUs \& DSPs:} dsPIC33CK (16-bit DSC, Microchip HAL), PIC32 (32-bit MCU, FreeRTOS), TMS320 DSP (bare-metal C, TI HAL), STM32 (bare-metal C, STM32 HAL), MSP430FR2355 (bare-metal C)\\
\textbf{Communications:} CAN 2.0/CAN FD; J1939, DroneCAN; I2C, SPI, UART/RS-232; IPMI/IPMB\\
\textbf{Systems, Debug \& Test:} TI AM62 Embedded Linux, Yocto/Arago, kernel and device trees, systemd, Docker; Python, Bash, Git/GitLab, JTAG, KiCad, SocketCAN/PCAN, oscilloscopes, logic/protocol analyzers, professional through-hole/SMT soldering

\section{Engineering Experience}
\companyheader{AEGIS Power Systems}
\roleentry{Embedded Firmware Engineer}{May 2026--Present}
\begin{resumebullets}
  \item Develop 16-bit dsPIC33CK firmware and bring up peripherals one subsystem at a time; isolate register, timing, configuration, and hardware faults with targeted tests, oscilloscopes, logic analyzers, schematics, and datasheets.
  \item Implemented and validated the VITA 46.11 Tier 2 management interface for a VPX power supply, providing IPMI 2.0 sensor, inventory, and event reporting over an I2C-based IPMB link to the chassis manager.
  \item Integrate battery hardware, CAN-connected devices, and Linux applications in embedded power systems; trace failures from power and wiring through bus traffic, services, configuration, and code.
  \item Support TI AM62 Embedded Linux integration: build Yocto/Arago images and kernels, update device trees, diagnose boot failures, and develop systemd services for external-watchdog recovery; use Docker for Windows/Linux tooling.
\end{resumebullets}

\roleentry{Embedded Systems Intern}{May--Aug 2025}
\roleentry{Computer Science Intern}{May--Aug 2024}
\descriptorentry{Selected work across both internships}
\begin{resumebullets}
  \item Developed firmware for a two-processor power system: a PIC32 running FreeRTOS managed I2C, SPI, CAN/J1939, and UART communications, while a TMS320 DSP ran bare-metal C firmware with TI hardware libraries for real-time power control.
  \item Developed DSP control firmware for transformer-coupled bidirectional buck--boost supplies using 12 linked PWM A/B pairs (24 outputs); configured each A channel's period/frequency, duty cycle, and relative phase while B followed or inverted A.
  \item Built host-side firmware-update tools for the same product: a C++ CLI accepted HEX files, while a Python TCP sender packetized updates with chunking and checksum validation; documented the update and toolchain workflows for engineering handoff.
  \item Implemented device-side update and recovery: the PIC32 staged firmware and a golden image in external SPI flash, programmed and verified the DSP over UART, and controlled the boot sequence.
  \item In 2024, created a restricted SSH/chroot interface with Bash and Python commands for secure remote inspection and control of TI Yocto/Arago hardware.
  \item Reverse-engineered RS-232 commands and checksums for poorly documented programmable resistive loads and relay controllers, then helped assemble the associated relay/load/power-supply test circuit.
\end{resumebullets}

\resumeentry{TEAM Industries | Engineering Intern}{Jun--Aug 2021 \& Jun--Aug 2022}
\begin{resumebullets}
  \item Developed a repeatable Excel/VBA dimensional and spline-tolerance validation tool that automatically flagged out-of-spec parts; built a searchable catalog of 17,000+ tooling, part, and vendor records with operation, location, and revision traceability.
\end{resumebullets}

\section{Senior Design}
\resumeentry{EV Active Sensor Adapter | ST MCU, TI BQ ADC, CAN}{NC State University}
\begin{resumebullets}
  \item Helped architect an EV sensor adapter integrating an STMicroelectronics MCU, TI BQ-family battery monitor/ADC, CAN transceivers, sensing, power, and telemetry on a custom PCB.
  \item Translated requirements and component trade studies into interface specifications and subsystem tests; coordinated design reviews, risk tracking, documentation, and bring-up across the team.
\end{resumebullets}

\section{Teaching \& Technical Leadership}
\resumeentry{ECE 306 Teaching Assistant \& Course Mentor | NC State University}{Spring 2025--Spring 2026}
\begin{resumebullets}
  \item Supported three semesters: two in an unofficial role at about 10 hr/week and one as the official TA at 20 hr/week; developed study guides and delivered test preparation, labs, code/PCB debugging, project guidance, and ECE-program guidance.
  \item Taught students to build and debug MSP430FR2355-based vehicle subsystems spanning ADC sensing, motor control, GPIO/interrupts, Wi-Fi/serial links, LCDs, and custom PCBs; added supplemental lectures and question-driven support as the class average rose from C+ the prior semester to A-.
\end{resumebullets}

\resumeentry{IEEE Open Project Space 2 | Lead Instructor}{Fall 2025 \& Spring 2026}
\begin{resumebullets}
  \item Led two semesters of a project-based electronics course from idea through demonstration; taught firmware, PCB/BOM workflows, and through-hole/SMT assembly while coaching students to solve their own design problems.
\end{resumebullets}

\resumeentry{IEEE Soldering Workshops | Coordinator, Lecturer \& Lab Instructor}{}
\begin{resumebullets}
  \item Led standalone workshops using custom through-hole and surface-mount flashing-light PCB keepsakes; taught professional soldering technique, safety, inspection, debugging, and rework.
\end{resumebullets}

\end{document}
